% =============================================================================
% PROPOSTA V2 DE CAPÍTULO REDUZIDO, FUNDAMENTAÇÃO TEÓRICA, SEÇÃO SEQUÊNCIAS DIDÁTICAS
% Versão: 2026-04-29 (segunda iteração)
% Estado anterior (v1): 14 a 15 páginas
% Estado proposto (v2): 13,5 a 14 páginas (redução adicional de ~1 página)
% Estratégia v2: aplicar Candidato 4 (remover SÍNTESE redundante de Mediação Docente)
%                e Candidato 5 (enxugar parágrafo BNCC com enumerações redundantes).
%                Núcleo teórico, identidade da sequência e subseções centrais
%                permanecem preservados.
% =============================================================================

\chapter{FUNDAMENTAÇÃO TEÓRICA}

{\color{blue} A fundamentação teórica representa a base epistemológica e metodológica desta investigação, pois explicita os referenciais que orientam tanto a análise do problema em estudo quanto a elaboração do artefato didático proposto. No âmbito da Educação Matemática, construir um referencial teórico vai além da simples apresentação de autores; envolve compreender de que modo distintas concepções de conhecimento, ensino e aprendizagem repercutem em decisões didáticas concretas (Artigue, 2014; Brousseau, 1997). No contexto desta pesquisa, esse referencial orienta a elaboração de uma sequência didática digital destinada ao ensino de Probabilidade no Ensino Médio, articulando princípios da Engenharia Didática, da Teoria das Situações Didáticas, dos ciclos iterativos de concepção em ambientes digitais e do construcionismo.

A opção pelo método Design Science Research (Dresch; Lacerda; Antunes Júnior, 2015) pressupõe que o artefato desenvolvido esteja ancorado em bases teóricas consistentes, uma vez que a intervenção pedagógica é concebida como resposta fundamentada a uma classe de problemas educacionais. Esta seção estrutura-se em três eixos: (i) as sequências didáticas como forma racional de organização do ensino; (ii) os fundamentos do ensino de Probabilidade no Ensino Médio; e (iii) o papel dos Objetos Virtuais de Aprendizagem na mediação do conhecimento matemático. O primeiro eixo é apresentado a seguir.

\section{Sequências Didáticas}

\hl{SEQUÊNCIA DIDÁTICA NA EDUCAÇÃO E NA MATEMÁTICA} A organização do ensino por meio de sequências didáticas firmou-se como uma das abordagens centrais na Educação Matemática contemporânea, contribuindo para o planejamento sistemático de situações orientadas por objetivos claros, progressão conceitual e articulação entre conteúdos, metodologias e avaliação. \citeonline{artigue2014} destaca que a elaboração de uma sequência requer análise epistemológica do conteúdo, antecipação de dificuldades cognitivas e controle interno das escolhas didáticas. \citeonline{brousseau2002}, por sua vez, argumenta que o ensino deve organizar situações que produzam necessidade intelectual real para o estudante, superando a lógica expositiva tradicional. Essa discussão torna-se particularmente relevante no ensino de conteúdos abstratos como a Probabilidade, cuja natureza contraintuitiva demanda organização progressiva que favoreça a construção de significados antes da formalização \cite[p.~9, 16, 23]{batanero2016}.

\hl{SEQUÊNCIA DIDÁTICA CONFORME ZABALA} \citeauthor{zabala1998} (\citeyear{zabala1998}, p.~18) define sequência didática como ``um conjunto de atividades ordenadas, estruturadas e articuladas para a realização de determinados objetivos educacionais, que têm um princípio e um fim conhecidos tanto pelo professor quanto pelos alunos'', evidenciando que ela possui unidade interna e direção formativa deliberada, e não apenas encadeamento de exercícios. No âmbito da Educação Matemática, esse conceito adquire densidade teórica própria, sendo ressignificado por tradições didáticas que o distanciam de uma compreensão estritamente operacional, em que cada etapa apresenta desafios conceituais que levam o estudante a rever e ampliar sua forma de pensar \cite[p.~66, 88]{brousseau2002}.

\subsection{Abordagens Metodológicas para a Construção de Sequências Didáticas}

\hl {PRINCÍPIOS E ABORDAGENS DA CONSTRUÇÃO DE UMA SEQUÊNCIA} A construção de uma sequência didática implica assumir o ensino como objeto de concepção racional, o que exige fundamentação metodológica explícita e critérios claros de validação. Saber o que uma sequência bem construída deve conter não responde à pergunta de como concebê-la, analisá-la e validá-la com rigor; essa questão remete às abordagens metodológicas que fundamentam as decisões de concepção, da análise preliminar do saber matemático à avaliação das aprendizagens efetivamente produzidas. No contexto brasileiro, destacam-se a Engenharia Didática, a Aprendizagem Baseada em Problemas e os ciclos iterativos de concepção. A Aprendizagem Baseada em Problemas \cite{viana2020abp} foi descartada como referência central por não dialogar com a estrutura de situação adidática individual mediada por Objeto Virtual de Aprendizagem que orienta esta pesquisa. As demais abordagens sustentam o produto desenvolvido e são examinadas a seguir.

\subsection{Engenharia Didática}

\hl {ABORDAGEM DA ENGENHARIA DIDÁTICA} No campo da Educação Matemática, a Engenharia Didática constitui uma das principais referências metodológicas quando a sequência didática está vinculada a um projeto de pesquisa. Artigue (1988) caracteriza essa abordagem como uma metodologia que articula a produção de situações de ensino com a investigação sistemática dos fenômenos didáticos delas decorrentes. A noção de ``engenharia'' não remete à aplicação de modelos prontos, mas à construção deliberada de situações fundamentadas teoricamente e submetidas a controle interno, organizada em quatro fases articuladas.

\hl {ANÁLISE PRELIMINAR} O ponto de partida é a análise preliminar, que examina o saber matemático em três dimensões. \citeonline[p.~68]{almoulouCoutinho2008} destacam que essa etapa envolve análises de natureza epistemológica (gênese histórica do conceito e obstáculos conceituais), didática (práticas correntes de ensino, materiais e restrições institucionais) e cognitiva (conhecimentos prévios, concepções espontâneas e dificuldades recorrentes apontadas pela literatura). A articulação entre essas dimensões fundamenta decisões didáticas teoricamente informadas.

\hl {CONCEPÇÃO, EXPERIMENTAÇÃO E VALIDAÇÃO} As três fases seguintes encadeiam-se em ciclo. Na concepção e análise a priori, o pesquisador delimita um campo de possibilidades e formula hipóteses sobre raciocínios, representações intermediárias e condições necessárias para o avanço conceitual \cite{almoulouSilva2012}, justificando cada escolha didática antes da implementação. Na experimentação, o professor assume simultaneamente o papel de mediador e de observador das interações efetivamente produzidas, admitindo correções e ajustes em diálogo com a análise a posteriori \cite{almoulouCoutinho2008}. Na análise a posteriori e validação, que encerram o ciclo, a validação realiza-se internamente, pela confrontação entre análise a priori e a posteriori \cite{artigue1996}: a qualidade do design depende da consistência das previsões formuladas, e quando estas são frágeis ou genéricas, a sequência tende a apoiar-se apenas em impressões subjetivas de funcionamento.

\subsection{Teoria das Situações Didáticas}

\hl {TSD DE BROUSSEAU} A Teoria das Situações Didáticas (TSD), desenvolvida por Guy Brousseau e sistematizada em sua obra de referência \cite{brousseau2002}, é uma das principais contribuições da Didática da Matemática francesa. Diferentemente da Engenharia Didática, que funciona como metodologia de pesquisa, a TSD é uma teoria sobre como o ensino deve ser organizado: o estudante aprende Matemática quando é colocado diante de situações que tornam o conhecimento necessário, e não quando recebe explicações prontas.

\hl {MILIEU E SITUAÇÃO ADIDÁTICA} Conceito-chave da teoria, o \emph{milieu} é o ambiente com o qual o estudante interage durante a atividade. Brousseau (2002) explica que esse ambiente fornece os retornos necessários para que o estudante avance, corrija estratégias, reformule hipóteses e construa respostas sem depender a cada passo da intervenção do professor. Quando o milieu está bem construído, emerge a \emph{situação adidática}, em que o estudante age por razões matemáticas próprias da situação e não para satisfazer uma exigência externa do professor.

\hl {DEVOLUÇÃO E MEDIAÇÃO DOCENTE} Para que uma situação adidática funcione, o professor precisa realizar a \emph{devolução}: transferir ao estudante a responsabilidade pela situação e resistir à tentação de intervir antes da hora \cite[p.~31]{brousseau2002}. Esse movimento opõe-se ao contrato didático implícito em que o professor apresenta o modelo e o estudante o repete; rompê-lo é condição para que a aprendizagem ocorra. A mediação docente assume, então, papel decisivo: intervir em excesso compromete a devolução, intervir de menos transforma o impasse em estagnação \cite{brousseau2002}, e decidir quando, como e com que intensidade intervir constitui competência docente central na gestão de sequências didáticas \cite[p.~159]{artigue2014}.

\hl{LIMITE EPISTEMOLÓGICO E INSTITUCIONALIZAÇÃO} A presença de um problema inicial não garante, por si só, a emergência de necessidades epistemológicas: para que haja reorganização conceitual, os efeitos do meio devem ser suficientes para provocar, por si mesmos, as adaptações esperadas, sob pena de tornar a situação adidática ``totalmente incapaz de provocar qualquer aprendizagem'' \cite[p.~58, tradução nossa]{brousseau2002}. As situações centrais desta sequência foram concebidas para que os retornos do Objeto Virtual de Aprendizagem confrontem o estudante com evidências que o cálculo intuitivo não consegue acomodar. O ciclo da TSD encerra-se apenas com a institucionalização, momento em que o professor reconhece o saber construído e lhe confere significado matemático formal \cite[p.~56]{brousseau1997}, conectando a experiência vivida ao conhecimento que a escola precisa ensinar.

\subsection{Particularidades das abordagens com OVAs}

\hl {OVAs, CICLOS ITERATIVOS E SIMULAÇÕES} Além das abordagens anteriores, ciclos iterativos de concepção têm orientado o desenvolvimento de sequências didáticas que integram tecnologias digitais. Sequências mediadas por tecnologia precisam ser concebidas, testadas e ajustadas continuamente, em ciclos que, embora menos formalizados que a Engenharia Didática, exigem o mesmo compromisso com o pensamento matemático dos estudantes. \citeonline{borbaScucugliaGadanidis2014} destacam que ambientes digitais interativos permitem ao estudante experimentar, visualizar e reformular ideias em tempo real. No ensino de Probabilidade, esse potencial é particularmente decisivo, pois \citeonline{kataokaOliveiraSouza2011} sustentam que a abordagem frequentista no Ensino Básico depende da possibilidade de executar repetições em escala inviável manualmente, observando a estabilização da frequência relativa em direção à probabilidade teórica. A cada ciclo, as situações são avaliadas quanto à sua capacidade de gerar uma necessidade real de aprender, preservando o princípio central da TSD de que a tecnologia fortalece o milieu, mas não substitui o desafio intelectual que deve mover o estudante: um simulador só cumpre função didática se estiver inserido em situação que exija interpretação e tomada de decisão.

\subsection{Como elaborar uma Sequência Didática no Ensino de Matemática}

\hl {COORDENAÇÃO DE REGISTROS E PROGRESSÃO} Apresentadas as abordagens metodológicas, cabe examinar como a estrutura se concretiza no planejamento. Diferentemente de outras disciplinas, a Matemática requer organização progressiva de conceitos interdependentes cuja compreensão depende da articulação entre formas de representação. Duval (2003, p.~21) afirma que a compreensão matemática não se limita à manipulação simbólica, mas demanda coordenação entre registros: é a capacidade de converter representações entre registros distintos que constitui o principal indicador de compreensão genuína. Daí a necessidade de prever atividades que integrem linguagem verbal, representação gráfica, notação simbólica e análise tabular, especialmente em Probabilidade, em que conceitos como independência de eventos ou probabilidade condicional exigem reestruturação do espaço amostral, sem a qual o conhecimento permanece superficial.

\hl{ELEMENTOS ESTRUTURAIS E ESTRUTURA PADRÃO} A identificação dos elementos estruturais que organizam uma sequência didática é passo decisivo para distingui-la de uma simples sucessão de tarefas. Artigue (2014, p.~160) sustenta que a coerência de uma sequência se fundamenta em análises que consideram, simultaneamente, a progressão conceitual e as interações previstas entre estudantes, conhecimento e meio didático. Estudos de revisão reforçam essa compreensão: ao analisarem trabalhos defendidos entre 2015 e 2019, \citeonline[p.~10]{lopesetal2020} registram a recorrência de um arranjo composto por cinco momentos (questionário inicial, introdução, problematização, discussão e resolução das atividades, e questionário final), arranjo que antecipa, no plano metodológico, o percurso adotado nesta pesquisa, em que a avaliação diagnóstica precede a intervenção com Objetos Virtuais de Aprendizagem e o questionário final encerra o ciclo.

\subsection{Progressão conceitual}

\hl {PROGRESSÃO E RACIOCÍNIO PROBABILÍSTICO} A progressão dos conteúdos matemáticos não pode ser definida apenas pela ordem de apresentação em livros didáticos ou documentos curriculares. Artigue (2014) sustenta que ela deve fundamentar-se em análises que articulem a lógica interna do saber e a lógica da aprendizagem, considerando dimensões epistemológicas e cognitivas. No campo da Probabilidade, certos conceitos precisam ser estabilizados antes de qualquer formalização algébrica, em especial o estatuto da aleatoriedade e o significado do espaço amostral. Nessa direção, \citeonline[p.~25]{cazorlaKataokaSilva2010} indicam que o desenvolvimento do raciocínio probabilístico requer a compreensão de que nem todos os fenômenos são determinísticos e a articulação entre a concepção clássica, em que a probabilidade é calculada a priori sobre o espaço amostral, e a concepção frequentista, em que ela emerge da estabilização da frequência relativa em grande número de repetições. Antes de formalizar, os estudantes precisam distinguir eventos certos, impossíveis e prováveis, interpretar regularidades empíricas e compreender que a frequência observada em poucas repetições não constitui uma ``lei do acaso''.

\hl {ORGANIZAÇÃO PROGRESSIVA E GÊNESE INSTRUMENTAL} A organização progressiva das situações de aprendizagem operacionaliza-se, na sequência aqui apresentada, pela exibição progressiva de pistas após erros sucessivos do estudante e por mensagens de andaime do tipo \emph{Lembre-se} nos momentos de transição conceitual; tais regulações pertencem ao funcionamento do milieu \cite{brousseau2002} e à análise a priori da Engenharia Didática \cite{artigue2014}. As pesquisas em Educação Matemática reconhecem ainda o duplo processo pelo qual estudante e instrumento digital se moldam reciprocamente, a instrumentação e a instrumentalização \cite{trouche2004}, o que justifica a integração intencional dos recursos digitais à sequência, para além do suporte técnico.

\subsection{Problematização}

\hl{PROBLEMATIZAÇÃO LEGÍTIMA E CONHECIMENTO PRÉVIO} Do ponto de vista da aprendizagem significativa, \citeonline[p.~ix]{ausubel2000} afirma que o fator mais importante que influencia a aprendizagem é aquilo que o aluno já sabe, princípio operacionalizado nesta sequência pela retomada explícita do experimento aleatório do disco antes da introdução do dado. Esse ponto de partida exige cuidado porque os estudantes costumam chegar à sala de aula com ideias equivocadas sobre acaso e aleatoriedade, o que \citeonline{brousseau2002} denomina obstáculos epistemológicos, que precisam ser estruturalmente enfrentados pela sequência. Brousseau \cite[p.~22-23, tradução nossa]{brousseau1997} sustenta que ``o saber deve aparecer como solução para um problema que o aluno reconhece como seu'', deslocando o papel do professor de transmissor de respostas para organizador de situações que provoquem desequilíbrio cognitivo produtivo. Nem todo problema proposto constitui, contudo, uma problematização no sentido epistemológico: em muitos casos, o que se denomina ``problema'' é apenas exercício contextualizado cuja resolução já está implicitamente definida pelo modelo previamente apresentado.

\hl {TRÍADE OBJETIVOS, ATIVIDADES E AVALIAÇÃO} Reconhecer essa distinção é apenas o primeiro passo: mesmo uma problematização legítima não sustenta, por si só, a aprendizagem pretendida. Como observa Zabala (1998, p.~38, p.~42), a aprendizagem só se torna consistente quando o problema está articulado a objetivos claros, atividades bem escolhidas e critérios de avaliação coerentes, e as atividades devem organizar-se de modo a permitir reconstrução progressiva dos significados, retomando o problema inicial ao longo da sequência. Em termos práticos, isso implica formular objetivos conceitualmente precisos, como identificar corretamente o espaço amostral em experimentos compostos ou interpretar a probabilidade condicional como redefinição do espaço amostral, em vez de formulações genéricas. Batanero e Díaz (2007, p.~128, tradução nossa) observam que ``os estudantes frequentemente recorrem a heurísticas intuitivas que entram em conflito com o raciocínio probabilístico normativo'', indicando que os objetivos da sequência devem incluir a superação de vieses cognitivos persistentes.

\subsection{Coerência}

\hl {COERÊNCIA INTERNA E AVALIAÇÃO} A coerência interna manifesta-se na articulação entre objetivos de aprendizagem, atividades propostas e critérios de avaliação. Zabala (1998) destaca que as atividades não são neutras: incorporam concepções implícitas sobre o conhecimento e o papel atribuído ao estudante. Quando os objetivos privilegiam a construção conceitual, as atividades devem criar situações em que os conceitos se tornem necessários para a resolução de problemas significativos; quando envolvem procedimentos, devem oferecer oportunidades de uso e refinamento em contextos progressivamente mais complexos; quando se almeja o desenvolvimento de atitudes como argumentação e tomada de decisão sob incerteza, devem favorecer a reflexão e o confronto de estratégias. A avaliação, por sua vez, não pode ser concebida como etapa externa: \citeonline[p.~63]{zabala1998} argumenta que ela deve estar integrada à sequência, funcionando como instrumento regulador da aprendizagem, e demandar argumentação e tomada de decisão em vez de aplicação mecânica de procedimentos.

\hl {O ERRO E A ARTICULAÇÃO DE REGISTROS} O erro assume papel estruturante, pois decorre da lógica interna do meio e não de correções externas que substituam a reflexão do estudante \cite[p.~88]{brousseau1997}. O funcionamento didático pressupõe situações que admitam tentativas iniciais e retornos sucessivos ao problema, organizando situações que se repetem, se transformam e evoluem, preservando uma situação-base comum, sem reduzir o conhecimento a procedimentos descontextualizados \cite[p.~31, p.~231]{brousseau1997}. É nessa recorrência controlada que o estudante reconhece regularidades, distingue invariantes e atribui sentido às representações formais. A coerência estrutural exige também articulação entre registros de representação, princípio já apresentado a partir de Duval (2003, p.~21), o que, no ensino de Probabilidade, significa integrar tabelas de contingência, diagramas de árvore, diagramas de Venn, linguagem verbal e expressão algébrica.

Compreendidos os elementos da sequência, cabe examinar o papel do docente nesse contexto.

\subsection{Papel da mediação docente na aprendizagem}

\hl {DEVOLUÇÃO E TENSÃO DIDÁTICA} Entre a aceitação do problema e a produção de respostas, a docência caracteriza-se por uma contenção deliberada: o conhecimento deve emergir do funcionamento da situação e das exigências de justificação que ela impõe. Cada conteúdo matemático pode ser associado a uma ou mais situações adidáticas fundamentais, cujo arranjo, orientado por finalidades de ensino, define tanto o saber mobilizado quanto o significado que ele assume \cite[p.~30-31, p.~230]{brousseau1997}. Essa contenção, contudo, gera tensão didática que não pode ser ignorada: ao transferir ao estudante a responsabilidade pela situação adidática, cria-se condição para uma atividade intelectual autêntica, mas também risco de impasses prolongados, dispersão ou soluções de curto alcance. O planejamento da sequência precisa antecipar esse risco e definir critérios de observação do progresso conceitual, evitando que silêncio ou hesitação sejam confundidos com aprendizagem.

\hl {GESTÃO, FORMALIZAÇÃO E CONSTRUCIONISMO} Cabe também ao professor a gestão do tempo didático e da memória do sistema de ensino, decidindo quando determinados conhecimentos devem ser retomados, transformados ou temporariamente suspensos. A mediação envolve duas responsabilidades complementares: criar condições para que o estudante construa conhecimento e garantir que esse conhecimento seja reconhecido e formalizado. Quando essa construção envolve a edificação progressiva de um artefato matemático com significado pessoal, como ocorre no laboratório de diagramas de Venn da sequência aqui apresentada, opera-se também o construcionismo de \citeonline{papert1994}, segundo o qual o aprender se faz pelo fazer com tecnologia. Essa articulação entre ação do professor, conteúdos probabilísticos e exigências curriculares é examinada na subseção seguinte.

\subsection{Sequências Didáticas para o Ensino de Probabilidade: articulações com a BNCC e recursos digitais}

\hl{BNCC E SEQUÊNCIA DIDÁTICA EM PROBABILIDADE} A Base Nacional Comum Curricular \cite[p.~520, 523, 528, 531]{brasil2018} destaca que o ensino de Probabilidade deve promover a compreensão de fenômenos aleatórios, o raciocínio em contextos de incerteza e a tomada de decisões fundamentadas, prevendo no Ensino Médio o uso de simulações e situações do mundo real. Tais competências não se desenvolvem por meio de ensino fragmentado, mas requerem planejamento sistemático e progressivo, articulado às competências da BNCC, com contextualização significativa e uso intencional de recursos digitais. O problema central reside em integrar tais recursos ao núcleo matemático da situação didática, evitando que a tecnologia se limite a um adorno metodológico, e em garantir que problemas relacionados a jogos, riscos, decisões cotidianas ou interpretação de dados estatísticos não sejam reduzidos à aplicação automática de fórmulas.

\hl {INTERATIVIDADE E INTENCIONALIDADE PEDAGÓGICA} Os recursos tecnológicos contemporâneos ampliam as possibilidades de trabalho com experimentos aleatórios. Ambientes como o GeoGebra possibilitam a manipulação de variáveis, a visualização dinâmica de espaços amostrais e a testagem de conjecturas. Ainda assim, a interatividade, por si só, não garante construção conceitual: um recurso digital pode acelerar respostas e reduzir reflexão se o design das tarefas não exigir justificativa, comparação de explicações e tomada de posição diante de resultados contraditórios. O meio tecnológico precisa, portanto, ser concebido como parte integrante da situação didática, com intencionalidade pedagógica clara, em que simulações sustentam confrontos produtivos entre intuição e evidência empírica e representações dinâmicas apoiam a formalização gradual dos conceitos \cite{trouche2004}. O ganho efetivo aparece quando o percurso inclui debate coletivo, sistematização conceitual e validação de argumentos, permitindo ao estudante compreender por que determinado conceito se torna necessário e em que condições ele se aplica.

\hl {IDENTIDADE DA SEQUÊNCIA DIDÁTICA DESTA PESQUISA} O rigor do percurso reside na coerência das escolhas didáticas e na capacidade de o conjunto das situações sustentar a reorganização conceitual do estudante, em vez de apenas na quantidade de atividades. A partir dos fundamentos discutidos, a sequência didática aqui proposta caracteriza-se como uma sequência didática digital, estruturada metodologicamente pela Engenharia Didática \cite{artigue2014} e operacionalizada por situações adidáticas \cite{brousseau2002} mediadas por dois Objetos Virtuais de Aprendizagem autorais para o ensino de Probabilidade no Ensino Médio. A Engenharia Didática orienta a concepção em quatro fases articuladas, da análise preliminar à validação interna; a Teoria das Situações Didáticas define o funcionamento do milieu computacional, no qual os retornos do artefato confrontam o estudante com evidências que sustentam a reorganização conceitual; e os Objetos Virtuais de Aprendizagem assumem a dupla função de meio interativo, em que se dá a devolução do problema, e de instrumento, no sentido da gênese instrumental de \citeonline{trouche2004}, com momentos pontuais de construcionismo \cite{papert1994} quando o estudante edifica artefatos matemáticos com significado pessoal. Essa articulação configura o tipo específico de sequência didática que esta pesquisa desenvolve, valida e analisa.
